\begin{flushright}
  \fechaclase{11 de noviembre de 2026}
\end{flushright}

con 

\begin{equation*}
    \begin{aligned}
        \epsilon_1=e_1^*-e_1\\
        \epsilon_2=e_2^*-e_2
    \end{aligned}
\end{equation*}

derivando

\begin{equation*}
    \begin{aligned}
        \dot{\epsilon}_1=\dot{e}_1^*-\dot{e}_1=e_2^*-e_2=\epsilon_2\\
        \dot{\epsilon}_2=\ddot{e}_1^*-\dot{e}_2=\ddot{e}_1^*-\ddot{r}+f(x,t+u)\\
    \end{aligned}
\end{equation*}

con concdiciones iniciales

\begin{equation*}
    \begin{aligned}
        \epsilon_1(0)=e_1^*(0)-e_1(0)=a-a=0\\
        \epsilon_2(0)=e_2^*(0)-e_2(0)=b-b=0
    \end{aligned}
\end{equation*}

se propone la superficie de deslizamiento

\begin{equation*}
    s=\epsilon_2+m\epsilon_1
\end{equation*}

con derivada

\begin{equation*}
    \begin{aligned}
        \dot{s}=&\dot{\epsilon}_2+m\dot{\epsilon}_1=\dot{\epsilon}_2+m\epsilon_2\\
            =& \ddot{e}_1^*-\ddot{r}+f(x,t)+m\epsilon_2+u
    \end{aligned}
\end{equation*}

proponemos la funcion de Lyapunov $V=\frac{1}{2}s^2$ y derivamos

\begin{equation*}
    \dot{V}=s\dot{s}=s(\ddot{e}_1^*-\ddot{r}+f(x,t)+m\epsilon_2+u)
\end{equation*}

asumiendo que $|-\ddot{r}+f(x,t)|\leq F$ para algun $F>0$

\begin{equation*}
    \dot{V} \leq s(\ddot{e}_1^*+F+m\epsilon_2+u)
\end{equation*}

proponemos el controlador

\begin{equation*}
    u=-\ddot{e}_1^*-m\epsilon_2-K Sign(s)
\end{equation*}

y sustituyendo

\begin{equation*}
    \begin{aligned}
        \dot{V}\leq& F|s|-K|s|\\
            \leq & -(K-F)|s|
    \end{aligned}
\end{equation*}

mostrando que al elegir $K>F$ el sistema es asintoticamente estable con tiempo de alcance

\begin{equation*}
    tr \leq \frac{|s(0)|}{K-F}
\end{equation*}

que debe cumplot $tr << T$. Finalmente el controlador SMC-PTC es dado por:

\begin{equation*}
    \begin{aligned}
        e_1=&r-y\\
        e_2=&\dot{r}-x_2\\
        a=& e_1(0)\\
        b=& e_2(0)\\
        c=& 2a+bT\\
        \tau = & \frac{t}{T}\\
    \end{aligned}
\end{equation*}

\begin{equation*}
    e_1^* = \begin{cases}
			(1-\tau)^2(a+c\tau), & 0 \leq t < T\\
            0, & t \geq T
		 \end{cases}
\end{equation*}

\begin{equation*}
    e_2^* = \begin{cases}
			b-\frac{2}{T}(3a+2bT)\tau+\frac{3c}{T}\tau, & 0 \leq t < T\\
            0, & t \geq T
		 \end{cases}
\end{equation*}

\begin{equation*}
    e_3^* = \begin{cases}
			-\frac{2}{T^2}(3a+2bT)+\frac{6c}{T^2}\tau, & 0 \leq t < T\\
            0, & t \geq T
		 \end{cases}
\end{equation*}

Por lo que el contorlador es

\begin{equation*}
    \begin{aligned}
        \epsilon_1=&e_1^*-e_1\\
        \epsilon_2=&e_2^*-e_2\\
        s = & \epsilon_2+m\epsilon_1\\
        u = & -\ddot{e}_1^*-m\epsilon_2-KSign(s)
    \end{aligned}
\end{equation*}

\textbf{Ejemplo:} Diseñe un SMC\_PTC para seguimiento de trayectorias para el sistema

\begin{equation*}
    \begin{aligned}
        \dot{x}_1=&x_2\\
        \dot{x}_2=&sin(x_1+x_2)+u\\
        y=&x_1
    \end{aligned}
\end{equation*}